\documentclass[a4paper]{article}
\usepackage{MMVII}

\newcommand{\TheCmd}{OriPoseEstimFocale}

\begin{document}

\CmdTitle
%\maketitle

\tableofcontents

 %=================================================================================
\section{Precaution}

{\center{\textcolor{red}{\bf ============== WARNING ===========}}}

Because of how an MMVII project is structured, the focal length and camera model information must be filled 
in for the command to work. Of course, since the goal is to estimate their values from scratch, 
the actual values entered do not matter. It is even recommended, if anything, to enter obviously absurd values, 
which will then be overwritten.

An example command with the “niche” dataset could be:

{\small
\begin{verbatim}
MMVII EditCalcMTDI Focalmm Save=1 Value=1e-3
MMVII EditCalcMTDI ModelCam PatternApply=.* Save=true Value=UnknownCam

MMVII OriPoseEstimFocale "IMG_569[456].JPG" TMP  InTieP=V1 Proj=FE_EquiDist
\end{verbatim}
}



 %=================================================================================

\SecDescr

In the standard photogrammetric pipeline, an approximate value of the camera parameters 
is initially required. For all parameters except the focal length, this approximate value
 can be obtained without any external information. The principal point can be placed at the 
center of the image (except for the rare tilt-shift or decentering lenses). For distortion, 
it can be assumed to be small, so that a zero value is a good approximation (provided that 
the lens type is known: rectilinear, equidistant fisheye, etc.). The focal length is a different 
story, and external information is generally needed: knowledge of the lens or metadata.


This command finds an initial focal length value from a multi-stereoscopic 
acquisition in which tie points have been extracted.
The value found is simply written to the console.  At the end, the following values
are printed for focal evaluation : 

\begin{itemize}
  \item the $35mm$ equivalent 
  \item  the focal in pixel;
  \item  if the camera is known, the "real" focal in mm;
\end{itemize}

If instead of tie points, we have ground control point, we can use the command
\FileSty{OriPoseEstim11P}.

 %=================================================================================

\section{Methodology}

The methodology is fairly basic: different focal lengths are tested, and the one yielding the best residuals 
is selected. For good results, at least three images are needed, containing enough triple points and having 
a sufficient stereoscopic baseline (say 10\%). For reasons that would take too long to detail here, 
if there are only pairwise tie points, the computation is very unstable and the results are generally unusable. 
That said, in the experiments we have carried out, increasing the number of images does not bring much improvement 
(although all the images provided will be used in the computation, at the cost of a longer computation time).

To perform this estimation, the command \FileSty{OriPoseEstim11P} is called for each focal length.


The focal lengths are chosen following a geometric progression. After computing the “best focal length”, 
if it is not one of the extreme values, the two neighboring values are used to fit a parabola and thereby 
refine the solution.


 %=================================================================================

\section{Mandatory and optional parameters}

%  -  -  -  -  -  -  -  -  -  -  -  -  -  -  -  -  -  -  -  -  -  -  -  -  -

\subsection{Mandatory parameters}

The two mandatory parameters are:

\begin{itemize}
  \item the set of images to which it applies;
  \item a location for  relative orientation that will be used by \FileSty{OriPoseEstim11P};
\end{itemize}

%  -  -  -  -  -  -  -  -  -  -  -  -  -  -  -  -  -  -  -  -  -  -  -  -  -

\subsection{Optional parameters}

The optionnal parameters are:

\begin{itemize}
   \item \FileSty{ InTieP, InMulTieP, InObjMesInstr} for tie points extraction (like \FileSty{OriPoseEstimRelGlob})
   \item \FileSty{Proj}  indicate model of camera
   \item \FileSty{IntFoc35}  indicates the interval of focal in $35mm$ equivalent; default value are very large :
          $[15,800]$ for rectilinear, $[6,25]$ for all fishe eyes and $[3,15]$ for equirectangular;
   \item \FileSty{RatioFocal}  indicate the ratio of geometric progression
   \item \FileSty{NbStep}  indicate a fix number of step instead of ratio
\end{itemize}

 %=================================================================================

\SecRelatedCmd


\begin{itemize}
   \item  \FileSty{OriPoseEstim11P}, does related things when we have GCP;
   \item  \FileSty{EditCalcMTDI} can be use before the command;
   \item  \FileSty{OriPoseEstimRelGlob} is used by this command.
\end{itemize}


\end{document}

